\begin{figure}
\begin{lstlisting}[
    commentstyle=\small\ttfamily\textcolor{gray},
    basicstyle=\small\ttfamily,
    xleftmargin=2.5em,
    numbers=left]
# Original term $t$, set of equational rewrite rules $\ruleset_\semeq$, maximum library size $N$
def LLMT($t$, $R_\semeq$, $N$):
    # EqSat phase:
    $\mathcal{G}$ = egraph($t$)                            # initialize with a single term $t$
    $\mathcal{G}$ = eqSat($\mathcal{G}, \ruleset_\semeq$)  # $\mathcal{G}$ represents all terms that are $\semeq t$

    # candidate generation phase:
    $\patset$ = AU($\mathcal{G}$)                          # generate candidate patterns by anti-unification
    $\ruleset_\kappa$ = $\{\kappa(p) \mid p \in \patset\}$ # construct a compression rule from every pattern
    $\mathcal{G}'$ = eqSat($\mathcal{G}, \ruleset_\kappa$) # $\mathcal{G}'$ represents all ways to compress $\mathcal{G}$

    # candidate selection phase:
    $\ruleset'$ = select_library($\mathcal{G}', N$)        # select the best $N$ rules from $\ruleset_\kappa$ using beam search
    $\mathcal{G}''$  = eqSat($\mathcal{G}, R'$)            # $\mathcal{G}''$ represents all ways to compress $\mathcal{G}$ using the optimal library          
    return extract($\mathcal{G}''$)                        # extract the smallest compressed term from $\mathcal{G}''$
\end{lstlisting}
  \caption{The top-level LLMT algorithm.}
\label{fig:algo-llmt}
\end{figure}

%\begin{figure}
%\begin{lstlisting}[
%    basicstyle=\normalsize\ttfamily,
%    xleftmargin=2.5em,
%    numbers=left]
%def anti-unify:
%    input: E: an e-graph containing the set of programs to anti-unify
%    output: R: a set of rewrites which introduces function definitions and applications
%               at the point of use
%
%    # TODO: i forgot how this works lmao
%\end{lstlisting}
%\caption{The antiunification algorithm of \babble.}
%\label{fig:algo-au}
%\end{figure}
